\documentclass[11pt,a4paper]{article}
\usepackage[T1]{fontenc}
\usepackage[utf8]{inputenc}
\usepackage{amsmath,amssymb,amsthm}
\usepackage[margin=1in]{geometry}
\usepackage[colorlinks=true,linkcolor=blue,urlcolor=blue]{hyperref}
\theoremstyle{plain}
\newtheorem{theorem}{Theorem}
\newtheorem{lemma}[theorem]{Lemma}
\newtheorem{proposition}[theorem]{Proposition}
\newtheorem{corollary}[theorem]{Corollary}
\theoremstyle{definition}
\newtheorem{remark}[theorem]{Remark}

\title{Recovering the unwritten recurrences of the table OEIS A237167}
\author{Adrian Perez Fontelles\\ \small Independent researcher}
\date{2 September 2026}
\begin{document}
\maketitle

\begin{abstract}
OEIS A237167 is a two-dimensional table $T(n,k)$ whose empirical recurrences are, for several of
its lines, given only as an ORDER --- ``k=2: [order 16]'' and the like --- with the
coefficients in a linked file rather than in the entry. Those conjectures can still be settled,
and the recurrences recovered. A line of the table is a fixed-width or fixed-height array
count, hence a walk count on finitely many vertices, hence satisfies some monic recurrence of
order at most that number; Berlekamp--Massey on twice as many exact terms returns the MINIMAL
one. Where its order equals the order the entry states --- and where the same holds after the
first few terms are dropped --- any recurrence of that order the line satisfies has a
characteristic polynomial that is a multiple of the minimal one and of equal degree, hence is
that polynomial. This note settles column 2, column 3, column 4.
\end{abstract}

\noindent\small 2020 Mathematics Subject Classification. 05A15, 11B37, 15A18.\normalsize

\section{The table and the unwritten conjectures}

OEIS A237167 is ``T(n,k)=Number of (n+1)X(k+1) 0..2 arrays with the maximum plus the minimum minus the upper median minus the lower median of every 2X2 subblock equal.''. It is read by antidiagonals and begins
\[
81, 245, 245, 777, 997, 777, 2597, 4365,\ \dots
\]
The entry carries, among its empirical claims:
\begin{quote}\small
k=2: [order 16]\\
k=3: [order 30]\\
k=4: [order 59]
\end{quote}
As of the ``Last modified'' line on the live entry (Jul 23 2025 09:55:28, revision 6) these are still
recorded as empirical, and the recurrences themselves are not in the entry text.

\section{Why a line of the table is a walk count}

Fix the index that the line fixes. The entry defines $T(n,k)$ as a number of arrays of a shape
determined by $n$ and $k$, subject to a condition local to a bounded window of consecutive
lines. Substituting the fixed index into the entry's own wording turns the definition into an
ordinary one-parameter array count, and for such a count the admissible windows are the
vertices of a finite digraph and an array is a walk. So the line is a walk count on $S$
vertices, and by the Cayley--Hamilton theorem it satisfies a monic integer recurrence of order
at most $S$.

\section{Recovering the recurrences}

\begin{lemma}\label{lem:bm}
A sequence known to satisfy some linear recurrence of order at most $S$ is determined, with its
minimal recurrence, by its first $2S$ terms; Berlekamp--Massey applied to them returns that
minimal recurrence exactly.
\end{lemma}

\subsection*{Column $k=2$}
Substituting the fixed index into the entry's wording gives an ordinary one-parameter array
count, a walk on $S=81$ vertices. Berlekamp--Massey on $2S=162$ exact terms
returns a minimal recurrence of order $16$ --- the order the entry states --- with
integer coefficients
\begin{center}\small\begin{tabular}{cccccccc}
$c_{1}$ & $15$ & $c_{5}$ & $391$ & $c_{9}$ & $-39520$ & $c_{13}$ & $-67104$ \\
$c_{2}$ & $-49$ & $c_{6}$ & $-13821$ & $c_{10}$ & $-137336$ & $c_{14}$ & $-119232$ \\
$c_{3}$ & $-250$ & $c_{7}$ & $7384$ & $c_{11}$ & $78264$ & $c_{15}$ & $20736$ \\
$c_{4}$ & $1567$ & $c_{8}$ & $58596$ & $c_{12}$ & $179256$ & $c_{16}$ & $31104$ \\
\end{tabular}\end{center}
so that $a(m)=\sum_{i=1}^{16}c_i\,a(m-i)$ for every $m$. Removing the first
$16+4$ terms and repeating gives order $16$ again, which is what the
identification below needs.

\subsection*{Column $k=3$}
Substituting the fixed index into the entry's wording gives an ordinary one-parameter array
count, a walk on $S=209$ vertices. Berlekamp--Massey on $2S=418$ exact terms
returns a minimal recurrence of order $30$ --- the order the entry states --- with
integer coefficients
\[
(c_1,\dots,c_{30})=(22,\ -99,\ -960,\ 8098,\ 12032,\ -218957,\ 19862,\ 3162005,\ -1747808,\ -28052026,\ 17465512,\ 161713315,\ -88019074,\ -621523196,\ 264308654,\ 1611400334,\ -500177858,\ -2823868897,\ 607262224,\ 3318019822,\ -470911456,\ -2561669588,\ 227604704,\ 1252643832,\ -64982496,\ -364185088,\ 9789184,\ 55881600,\ -579072,\ -3354624).
\]
so that $a(m)=\sum_{i=1}^{30}c_i\,a(m-i)$ for every $m$. Removing the first
$30+4$ terms and repeating gives order $30$ again, which is what the
identification below needs.

\subsection*{Column $k=4$}
Substituting the fixed index into the entry's wording gives an ordinary one-parameter array
count, a walk on $S=549$ vertices. Berlekamp--Massey on $2S=1098$ exact terms
returns a minimal recurrence of order $59$ --- the order the entry states --- with
integer coefficients
\[
(c_1,\dots,c_{59})=(38,\ -371,\ -2838,\ 64993,\ -58567,\ -4394237,\ 15058024,\ 172006601,\ -838703867,\ -4505621654,\ 26693317114,\ 85503737174,\ -572478983494,\ -1237303411668,\ 8868339925003,\ 14117658629076,\ -103158791332900,\ -129518825824655,\ 924072837087924,\ 964103900257693,\ -6489848726510141,\ -5838242601572712,\ 36216416422429550,\ 28752450595327795,\ -162225277010758868,\ -115017716694513491,\ 587691474722981835,\ 373164028715470753,\ -1731047885680265093,\ -980075261566163264,\ 4159412279471724748,\ 2078051147660241412,\ -8164286576577101504,\ -3542304485928678000,\ 13084030130740708600,\ 4824310506366988372,\ -17077539997844469364,\ -5201108564881985976,\ 18071183306742752088,\ 4378777606691183872,\ -15397149068880798752,\ -2819908771789766592,\ 10462825501680041280,\ 1343176554987231904,\ -5598704889450166176,\ -444086661136863232,\ 2319888053356117248,\ 86410720616830464,\ -727960690761887232,\ -2463057123006720,\ 167819902964087040,\ -3533160413067264,\ -27223090020771840,\ 946858535276544,\ 2907812839624704,\ -106061071417344,\ -181780813578240,\ 4550416662528,\ 5003075911680).
\]
so that $a(m)=\sum_{i=1}^{59}c_i\,a(m-i)$ for every $m$. Removing the first
$59+4$ terms and repeating gives order $59$ again, which is what the
identification below needs.

\begin{theorem}
For each line above, the displayed recurrence holds for every index, and it is the recurrence
the entry records.
\end{theorem}

\begin{proof}
It holds by Lemma~\ref{lem:bm}. For the identification, the entry asserts that the line
satisfies SOME recurrence of the stated order, possibly only from some index on. Let $q$ be its
characteristic polynomial and $q_{\min}$ the minimal polynomial of the tail. Then
$q_{\min}\mid q$, and the second Berlekamp--Massey run --- on the line with its first few
terms removed --- shows $\deg q_{\min}$ equals the stated order, which is $\deg q$. Both are
monic, so $q=q_{\min}$.
\end{proof}

\section{Verification}

Three checks.

First, each line's model was compared against the entry's own published values for that line,
taken from the antidiagonal DATA read in either orientation and from the ``Table starts''
block, and agrees at every available term. This is the only point at which reading the entry's
English enters, and a misreading gives different counts.

Second, each recovered recurrence was evaluated directly on the model's values, with no
Berlekamp--Massey involved, and holds at every index tested.

Third, each order was computed twice by independent routes --- modulo a $61$-bit prime, where
every intermediate is a machine word, and again in exact rational arithmetic --- and once more
on the line with its first few terms dropped, which is what the identification requires. All
agree.

\begin{thebibliography}{9}
\bibitem{oeis} The OEIS Foundation, \emph{The On-Line Encyclopedia of Integer
Sequences}, \url{https://oeis.org}, 2026. Sequence A237167.
\bibitem{massey} J.~L.~Massey, Shift-register synthesis and BCH decoding, \emph{IEEE
Trans. Inform. Theory} \textbf{15} (1969), 122--127.
\bibitem{stanley} R.~P.~Stanley, \emph{Enumerative Combinatorics}, Volume 1, 2nd ed.,
Cambridge University Press, 2011. (Transfer-matrix method, Section 4.7.)
\bibitem{horn} R.~A.~Horn and C.~R.~Johnson, \emph{Matrix Analysis}, 2nd ed., Cambridge
University Press, 2013.
\end{thebibliography}

\end{document}
